\section{Horizon energy and half-action cell}
\label{rm:chap:horizonte-media-accion}

\subsection{Geometric conventions of the de Sitter patch}

Consider a four-dimensional de Sitter space with areal radius
\(R>0\), causal velocity \(c\), and gravitational constant \(G\). Adopt
the convention
\begin{equation}
 H_R=\frac cR,
 \qquad
 \Lambda_R=\frac3{R^2},
 \qquad
 \ell_P^2=\frac{G\hbar}{c^3}.
 \label{rm:eq:horizon-ds-conventions}
\end{equation}
The density of energy associated to the term cosmological is
\begin{equation}
 \varepsilon_\Lambda
 :=\frac{\Lambda_Rc^4}{8\pi G}.
 \label{rm:eq:horizon-vacuum-density}
\end{equation}
The spatially flat section of the patch contains, up to radius \(R\), the
volume
\begin{equation}
 V_R=\frac{4\pi}{3}R^3.
 \label{rm:eq:horizon-flat-volume}
\end{equation}
These conventions positively fix the quasilocal energy used in the
chapter. Choosing another sign for the de Sitter Killing charge
belongs to a distinct construction and must be translated before formulas
are compared.

\subsection{energy of vacuum and energy of Misner--Sharp}

In a spherically symmetric geometry, let
\begin{equation}
 \dd s^2=g_{ab}(x)\dd x^a\dd x^b
 +\mathcal R(x)^2\dd\Omega_2^2.
 \label{rm:eq:horizon-spherical-metric}
\end{equation}
The Misner--Sharp energy \cite{rm:MisnerSharp1964}, under the adopted positive convention, is
defined by
\begin{equation}
 E_{\rm MS}(\mathcal R)
 :=\frac{c^4\mathcal R}{2G}
 \left(1-g^{ab}\nabla_a\mathcal R\nabla_b\mathcal R\right).
 \label{rm:eq:horizon-Misner-Sharp}
\end{equation}
A marginal sphere satisfies
\begin{equation}
 g^{ab}\nabla_a\mathcal R\nabla_b\mathcal R=0.
 \label{rm:eq:horizon-marginal-sphere}
\end{equation}

\begin{theorem}[Confluence volumetrica and cuasilocal]
\label{rm:thm:horizon-vacuum-MS}
On the de Sitter horizon \(\mathcal R=R\), the integrated vacuum
energy and the Misner--Sharp energy coincide exactly:
\begin{equation}
 \boxed{
 E_\Lambda
 :=\varepsilon_\Lambda V_R
 =E_{\rm MS}(R)
 =\frac{c^4R}{2G}.}
 \label{rm:eq:horizon-E-Lambda}
\end{equation}
\end{theorem}

\begin{proof}
The equations
\cref{rm:eq:horizon-ds-conventions,rm:eq:horizon-vacuum-density,rm:eq:horizon-flat-volume}
give
\[
 \varepsilon_\Lambda V_R
 =\frac{3c^4}{8\pi GR^2}\frac{4\pi R^3}{3}
 =\frac{c^4R}{2G}.
\]
The marginal condition \cref{rm:eq:horizon-marginal-sphere} applied to
\cref{rm:eq:horizon-Misner-Sharp} yields the same result.
\end{proof}

The equality brings together two geometric publications. The first integrates a
volumetric density; the second evaluates the trapping defect of the
areal sphere. Both preserve the radius \(R\) and gravitational compliance
\(G/c^4\).

\subsection{Temperature, entropy and product of horizon}

The de Sitter surface gravity has magnitude \(c^2/R\). The Gibbons--Hawking
temperature \cite{rm:GibbonsHawking1977} and the Bekenstein--Hawking
entropy \cite{rm:Bekenstein1973,rm:Hawking1975} are
\begin{align}
 k_BT_H&=\frac{\hbar c}{2\pi R},
 \label{rm:eq:horizon-temperature}\\
 \frac{S_H}{k_B}
 &=\frac{A_H}{4\ell_P^2}
 =\frac{4\pi R^2}{4\ell_P^2}
 =\frac{\pi R^2}{\ell_P^2}.
 \label{rm:eq:horizon-entropy}
\end{align}

\begin{theorem}[de Sitter thermodynamic identity]
\label{rm:thm:horizon-TS}
For every radius \(R\) of the declared family,
\begin{equation}
 \boxed{
 T_HS_H=\frac{c^4R}{2G}=E_\Lambda.}
 \label{rm:eq:horizon-TS-E}
\end{equation}
\end{theorem}

\begin{proof}
Multiplying \cref{rm:eq:horizon-temperature,rm:eq:horizon-entropy} gives
\[
 T_HS_H
 =\frac{\hbar c}{2\pi Rk_B}
  \frac{k_B\pi R^2}{\ell_P^2}
 =\frac{\hbar cR}{2\ell_P^2}.
\]
The identity \(\ell_P^2=G\hbar/c^3\) turns the last member into
\(c^4R/(2G)\), equal to \cref{rm:eq:horizon-E-Lambda}.
\end{proof}

The factor \(1/2\) has a decomposition geometric transparent. The
areal factor \(4\pi\), the normalization holographic by four and the
periodicity thermal \((2\pi)^{-1}\) satisfy
\begin{equation}
 \frac{4\pi}{4}\frac1{2\pi}=\frac12.
 \label{rm:eq:horizon-half-factor}
\end{equation}
The half is therefore the joint contraction between areal quantization
and Euclidean periodicity.

\subsection{How complete of the horizon and Planck force}

We define the tension or fuerza of Planck and the complete quantum associated with the
radius by
\begin{equation}
 \mathscr F_P:=\frac{c^4}{G},
 \qquad
 \boxed{
 \mathcal Q_H(R):=2T_HS_H
 =\frac{c^4R}{G}=\mathscr F_P\,R.}
 \label{rm:eq:horizon-full-quantum}
\end{equation}
Quasilocal energy occupies half a cell of the complete quantum:
\begin{equation}
 \boxed{E_\Lambda=T_HS_H=\frac{\mathcal Q_H(R)}2.}
 \label{rm:eq:horizon-half-quantum-general}
\end{equation}

\begin{proposition}[Radial compliance]
\label{rm:prop:horizon-compliance}
The derivative of the complete quantum with respect to radius is the Planck force,
and the derivative of the quasilocal energy is half of it:
\begin{equation}
 \frac{\dd\mathcal Q_H}{\dd R}=\mathscr F_P,
 \qquad
 \frac{\dd E_\Lambda}{\dd R}=\frac{\mathscr F_P}{2}.
 \label{rm:eq:horizon-radial-derivatives}
\end{equation}
The compliance gravitatoria \(\mathscr C_P=G/c^4\) satisfies
\begin{equation}
 R=\mathscr C_P\mathcal Q_H.
 \label{rm:eq:horizon-compliance-law}
\end{equation}
\end{proposition}

\begin{proof}
The three identities are direct derivatives or inversions of
\cref{rm:eq:horizon-full-quantum}.
\end{proof}

The equation \cref{rm:eq:horizon-compliance-law} expresses the radial response:
the complete quantum applies tension, and geometry publishes a radius through
universal compliance. The product \(T_HS_H\) occupies the quasilocal half
of that response.

\subsection{Specialization to the closure chronological of 108 states}

The chronological closure of \cref{rm:ch:cosmo} determines
\begin{equation}
 R_{108}=\ell_P,
 \qquad
 t_P=\frac{\ell_P}{c},
 \qquad
 T_{108}=2\pi t_P,
 \label{rm:eq:horizon-CT108-scales}
\end{equation}
and the Planck quantum
\begin{equation}
 E_P=\frac{\hbar}{t_P}
 =k_B\Theta_{\rm clk}\log3.
 \label{rm:eq:horizon-Planck-trit}
\end{equation}

\begin{theorem}[Cell CT108 of half action]
\label{rm:thm:horizon-half-action}
In the specialization \(R=R_{108}=\ell_P\),
\begin{align}
 \mathcal Q_H(\ell_P)
 &=\frac{c^4\ell_P}{G}=E_P,
 \label{rm:eq:horizon-QH-EP}\\
 \boxed{
 E_\Lambda=T_HS_H
 =\frac{E_P}{2}
 =\frac{k_B\Theta_{\rm clk}\log3}{2},}
 \label{rm:eq:horizon-CT108-half-energy}\\
 \boxed{E_\Lambda t_P=\frac{\hbar}{2},}
 \qquad
 \boxed{E_\Lambda T_{108}=\frac h2.}
 \label{rm:eq:horizon-half-actions}
\end{align}
\end{theorem}

\begin{proof}
The first equality is obtained by substituting
\(\ell_P=\sqrt{G\hbar/c^3}\) and
\(E_P=\sqrt{\hbar c^5/G}\). The
\cref{rm:eq:horizon-half-quantum-general} produces
\(E_\Lambda=E_P/2\), and
\cref{rm:eq:horizon-Planck-trit} supplies the informational chart. Finally,
\[
 E_\Lambda t_P=\frac{E_Pt_P}{2}=\frac\hbar2
\]
and, as \(T_{108}=2\pi t_P\) and \(h=2\pi\hbar\),
\[
 E_\Lambda T_{108}
 =\frac{E_P}{2}(2\pi t_P)=\pi\hbar=\frac h2.
\]
\end{proof}

The first equality in \cref{rm:eq:horizon-half-actions} publishes half a reduced action
over Planck time; the second publishes half a complete action
over the CT108 period. Both express the same cell in the angular
and circular charts.

The temperature and entropy are also determined:
\begin{equation}
 k_BT_H=\frac{E_P}{2\pi},
 \qquad
 \boxed{\frac{S_H}{k_B}=\pi.}
 \label{rm:eq:horizon-CT108-TS}
\end{equation}
The temperature converts the frequency of horizon in energy; the
entropy converts the area of the sphere of Planck in information. Their
product recovers the same half-quantum.

\subsection{Effective horizon weight and hyperbolic exponential}

We define the effective thermodynamic weight
\begin{equation}
 \Omega_{\rm eff}(R):=\exp\!\left(\frac{S_H(R)}{k_B}\right).
 \label{rm:eq:horizon-effective-weight}
\end{equation}
In the CT108 set,
\begin{equation}
 \boxed{\Omega_{\rm eff}(\ell_P)=\ee^\pi.}
 \label{rm:eq:horizon-effective-e-pi}
\end{equation}
The weight \(\Omega_{\rm eff}\) is a positive quantity defined by
entropy. Its real value expresses effective multiplicity; an interpretation
as an integer cardinality of microstates requires an additional discretization.

The areal--volumetric incidence operator is
\begin{equation}
 F_{\rm av}=\begin{pmatrix}0&1\\1&1\end{pmatrix},
 \qquad
 J_\varphi:=\frac{2F_{\rm av}^2-3I}{\sqrt5}.
 \label{rm:eq:horizon-J-phi}
\end{equation}
A direct calculation gives
\begin{equation}
 J_\varphi^2=I,
 \qquad
 \Spec(J_\varphi)=\{+1,-1\}.
 \label{rm:eq:horizon-J-involution}
\end{equation}
The hyperbolic flow at closure time \(\pi\) is
\begin{equation}
 \mathcal E_\pi
 :=\ee^{\pi J_\varphi}
 =\cosh\pi\,I+\sinh\pi\,J_\varphi,
 \label{rm:eq:horizon-hyperbolic-flow}
\end{equation}
and possesses spectrum
\begin{equation}
 \boxed{
 \Spec(\mathcal E_\pi)=\{\ee^\pi,\ee^{-\pi}\},
 \qquad
 \rho(\mathcal E_\pi)=\ee^\pi.}
 \label{rm:eq:horizon-hyperbolic-spectrum}
\end{equation}

\begin{corollary}[Horizon--Perron scalar confluence]
\label{rm:cor:horizon-Perron-scalar}
In the CT108 closure,
\begin{equation}
 \boxed{
 \frac{S_H}{k_B}
 =\log\Omega_{\rm eff}
 =\log\rho(\mathcal E_\pi)
 =\pi.}
 \label{rm:eq:horizon-Perron-log}
\end{equation}
\end{corollary}

\begin{proof}
The equations
\cref{rm:eq:horizon-CT108-TS,rm:eq:horizon-effective-e-pi,rm:eq:horizon-hyperbolic-spectrum}
give the three members.
\end{proof}

The scalar identity is exact and preserves two distinct genealogies. The
horizon produces \(\ee^\pi\) by exponentiation of the areal entropy;
Perron dynamics produces it as the expansive eigenvalue of an
involution that is hyperbolic. The terminal equality fixes a boundary condition
for the intertwiner and preserves the distinction of domains.

\subsection{Consistency of the Gravitational Transducer with the Schwarzschild Horizon}

A Schwarzschild horizon with the same area radius satisfies
\begin{equation}
 R_H=\frac{2GE_H}{c^4},
 \qquad
 E_H=\frac{c^4R_H}{2G}.
 \label{rm:eq:horizon-Schwarzschild-energy}
\end{equation}
The entropy areal coincides with \cref{rm:eq:horizon-entropy}, whereas the
temperature is
\begin{equation}
 k_BT_{\rm Sch}=\frac{\hbar c}{4\pi R_H}
 =\frac12k_BT_H\big|_{R=R_H}.
 \label{rm:eq:horizon-Schwarzschild-temperature}
\end{equation}
Consequently,
\begin{equation}
 E_H=2T_{\rm Sch}S_H,
 \qquad
 T_H^{\rm dS}=2T_{\rm Sch}
 \label{rm:eq:horizon-Schwarzschild-Smarr}
\end{equation}
for the same radius. In \(R_H=\ell_P\),
\begin{equation}
 E_H=\frac{E_P}{2},
 \qquad
 T_{\rm Sch}S_H=\frac{E_P}{4}.
 \label{rm:eq:horizon-Schwarzschild-Planck}
\end{equation}
The contrast separates the two geometries: they share quasilocal energy,
area, and entropy, and their surface gravities produce temperatures
in the ratio two.

With \(\beta=(k_BT_{\rm Sch})^{-1}\), the action Euclidean satisfies
\begin{equation}
 \frac{I_E}{\hbar}
 =\frac{E_H}{k_BT_{\rm Sch}}-\frac{S_H}{k_B}
 =2\pi-\pi=\pi
 \label{rm:eq:horizon-Schwarzschild-action}
\end{equation}
at the Planck radius. The same \(\pi\) appears as reduced entropy,
reduced Euclidean action, hyperbolic-flow time, and the logarithm of the
effective weight.

\subsection{Horizontal Horizon–Perron Entangler: Spectral compatibility between boundary energy and return memory}

The coincidence
\(\Omega_{\rm eff}=\rho(\mathcal E_\pi)\) fixes a spectral datum and leaves
open the relation between the complete objects. Let
\(\mathcal A_H^{(m)}\) denote the algebra of horizon observables at nonadic level
\(m\), let \(\omega_H^{(m)}\) denote its state, and let
\(\mathcal P_\varphi^{(m)}\) denote the Perron transfer system.

\begin{definition}[Thermodynamic--dynamic intertwiner]
\label{rm:def:horizon-Perron-intertwiner}
An admissible intertwiner is a family of unital completely positive
linear maps (UCP)
\begin{equation}
 \mathfrak I_m:\mathcal A_H^{(m)}\longrightarrow
 \mathcal A_\varphi^{(m)}
 \label{rm:eq:horizon-intertwiner}
\end{equation}
which satisfies:
\begin{enumerate}[label=\roman*)]
 \item preservation of the state, in
addition to unitality and complete positivity already
incorporated into the UCP condition;
 \item compatibility with nonadic refinement;
 \item transport of the involution interior--exterior toward
       \(J_\varphi\mapsto-J_\varphi\);
 \item commutation between the horizon modular semigroup and the transfer
       semigroup on the declared subspace;
 \item equality of the spectral radius at time \(\pi\), given by
       \cref{rm:eq:horizon-Perron-log}.
\end{enumerate}
\end{definition}

The fifth condition is necessary and scalar; the first four provide the
algebraic content that a numerical coincidence does not contain. An
intertwiner constructed in this manner would identify the effective
thermodynamic multiplicity with a dynamics of expansion and contraction
while preserving orientation, state, and refinement.

\begin{proofstatusbox}
The identities
\(E_\Lambda=E_{\rm MS}=T_HS_H=c^4R/(2G)\), the definition
\(\mathcal Q_H=\mathscr F_P\,R\), the CT108 specialization, the two
half-action cells, and the equality
\(S_H/k_B=\log\rho(\ee^{\pi J_\varphi})=\pi\) are exact results under
the declared conventions and de Sitter realization. The identification
of \(\Omega_{\rm eff}\) with a microscopic cardinal and the intertwiner
\(\mathfrak I_m\) belong to a delimited physical interface.
\end{proofstatusbox}


\begin{falsifierbox}
Geometric closure fails if \(\Lambda_R\ne3/R^2\), if the sphere of
radius \(R\) fails to satisfy the marginal condition, or if the integrated energy differs from
\(E_{\rm MS}(R)\). Thermodynamic closure fails if the temperature or the
area law produces \(T_HS_H\ne c^4R/(2G)\). The CT108 specialization
fails if \(R_{108}\ne\ell_P\), \(T_{108}\ne2\pi t_P\), or
\(E_P\ne k_B\Theta_{\rm clk}\log3\). The Perron interface is refuted by
an obstruction of positivity, state, orientation, or refinement, even
when the eigenvalue \(\ee^\pi\) agrees scalarly.
\end{falsifierbox}
